\chapter{第三章标题}

\section{引言}

\section{提出的方法或者框架名称}

\subsection{子标题1}

\subsection{子标题2}

\subsection{子标题3}

\subsection{子标题4}

\subsection{子标题5}

\subsection{子标题6}

\section{实验与结果分析}

\subsection{数据集介绍}

\subsection{数据预处理}

\subsection{实验细节}

\subsection{子标题4}

表~\ref{tab:tabel1.2}~和表~\ref{tab:tabel1.3}~分别汇报了详细比较结果，

\begin{table}[htb]
  \caption{xxxxx性能比较}
  \centering
  \wuhao
  % 调整列间距以美化排版
  \setlength{\tabcolsep}{6mm}{
  \renewcommand{\arraystretch}{1.5}
  \begin{tabular}{cccc}
    \toprule
    任务 & 方法 & 指标1（\%）  & 指标2（\%） \\
    \midrule
    \multirow{6}{*}{任务1} 
    & 1   & $00.00^* \pm 00.00$ & $00.00 \pm 00.00$ \\
    & 2   & $00.00^* \pm 00.00$ & $00.00 \pm 00.00$ \\
    & 3   & $00.00^* \pm 00.00$ & $00.00 \pm 00.00$ \\
    & 4   & $00.00^* \pm 00.00$ & $00.00 \pm 00.00$ \\
    & 5   & $00.00^* \pm 00.00$ & $00.00 \pm 00.00$ \\
    & 本章方法     & $\textbf{00.00} \pm 00.00$ & $\textbf{00.00} \pm 00.00$ \\
    \cmidrule{1-4}
    \multirow{6}{*}{任务2} 
    & 1   & $00.00^* \pm 00.00$ & $00.00 \pm 00.00$ \\
    & 2   & $00.00^* \pm 00.00$ & $00.00 \pm 00.00$ \\
    & 3   & $00.00^* \pm 00.00$ & $00.00 \pm 00.00$ \\
    & 4   & $00.00^* \pm 00.00$ & $00.00 \pm 00.00$ \\
    & 5   & $00.00^* \pm 00.00$ & $00.00 \pm 00.00$ \\
    & 本章方法    & $\textbf{00.00} \pm 00.00$ & $\textbf{00.00} \pm 00.00$ \\
    \cmidrule{1-4}
    \multirow{6}{*}{任务3} 
    & 1   & $00.00^* \pm 00.00$ & $00.00 \pm 00.00$ \\
    & 2   & $00.00^* \pm 00.00$ & $00.00 \pm 00.00$ \\
    & 3   & $00.00^* \pm 00.00$ & $00.00 \pm 00.00$ \\
    & 4   & $00.00^* \pm 00.00$ & $00.00 \pm 00.00$ \\
    & 5   & $00.00^* \pm 00.00$ & $00.00 \pm 00.00$ \\
    & 本章方法     & $\textbf{00.00} \pm 00.00$ & $\textbf{00.00} \pm 00.00$ \\
    \bottomrule
  \end{tabular}}
  \\ \vspace{1ex} % 表格与注脚之间留一点间距

  \vspace{-1ex} % 微调间距
  \begin{flushleft}
    \footnotesize 注：“*” 表示与本章方法相比具有显著性差异（$p < 0.01$，配对 t 检验）。
  \end{flushleft}
  \label{tab:tabel1.2}
\end{table}

\begin{table}[htb]
  \caption{xxxxx性能比较}
  \centering
  \wuhao
  % 调整列间距以美化排版
  \setlength{\tabcolsep}{6mm}{
  \renewcommand{\arraystretch}{1.5}
  \begin{tabular}{cccc}
    \toprule
    任务 & 方法 & 指标1（\%）& 指标2（\%） \\
    \midrule
    \multirow{6}{*}{任务1} 
    & 1     & $00.00^* \pm 00.00$ & $00.00 \pm 00.00$ \\
    & 2     & $00.00^* \pm 00.00$ & $00.00 \pm 00.00$ \\
    & 3     & $00.00^* \pm 00.00$ & $00.63 \pm 00.00$ \\
    & 4     & $00.00^* \pm 00.00$ & $00.00 \pm 00.00$ \\
    & 5     & $00.00^* \pm 00.00$ & $00.00 \pm 00.00$ \\
    & 本章方法     & $\textbf{00.00} \pm 00.00$ & $\textbf{00.00} \pm 00.00$ \\
    \cmidrule{1-4}
    \multirow{6}{*}{任务2} 
    & 1     & $00.00^* \pm 00.00$ & $00.00 \pm 00.00$ \\
    & 2     & $00.00^* \pm 00.00$ & $00.00 \pm 00.00$ \\
    & 3     & $00.00^* \pm 00.00$ & $00.00 \pm 00.00$ \\
    & 4     & $00.00^* \pm 00.00$ & $00.00 \pm 00.00$ \\
    & 5     & $00.00^* \pm 00.00$ & $00.00 \pm 00.00$ \\
    & 本章方法     & $\textbf{00.00} \pm 00.00$ & $\textbf{00.00} \pm 00.00$ \\
    \cmidrule{1-4}
    \multirow{6}{*}{任务3} 
    & 1    & $99.16 \pm 01.58$   & $99.16 \pm 01.58$ \\
    & 2    & $92.74^* \pm 06.41$ & $92.75 \pm 06.39$ \\
    & 3    & $96.11^* \pm 01.84$ & $96.13 \pm 01.83$ \\
    & 4    & $97.22^* \pm 01.24$ & $97.22 \pm 01.24$ \\
    & 5    & $98.33^* \pm 01.04$ & $98.33 \pm 01.04$ \\
    & 本章方法     & $\textbf{00.00} \pm 00.00$ & $\textbf{00.00} \pm 00.00$ \\
    \bottomrule
  \end{tabular}}
  \\ \vspace{1ex} % 表格与注脚之间留一点间距

  \vspace{-1ex} % 微调间距
  \begin{flushleft}
    \footnotesize 注：“*” 表示与本章方法相比具有显著性差异（$p < 0.01$，配对 t 检验）。
  \end{flushleft}
  \label{tab:tabel1.3}
\end{table}

\subsection{子标题5}

\section{讨论与分析}

\subsection{消融实验}

\subsection{子标题2}

\subsection{子标题3}

\subsection{子标题4}

\subsection{子标题5}

\subsection{局限性与未来工作}

尽管所提框架取得了有价值的实验结果，仍存在以下几方面值得深入探讨的局限性。

\section{本章小结}


